\scanpage{58}
et si on se donne un hom de $X$ dans un \deleted{principal homog.} schéma \unsure{paraabélien|parabélien}) sous un schéma abélien $A$) un hom $A_Y\longrightarrow A$, via un élément de ${\mathrm{CK}}^1(A_Y,A^*)$. \inserted{Sur les} fibres \deleted{par fibre}\inserted{\unsure{géom|géom.}}, c'est donné par la construction déjà explicitée\deleted{\unsure{)}}.

2. \uline{Application à l'algébricité de certains homomorphismes de V.A. ou d'\unsure{espaces} de cohomologie $l$-\unsure{adiques|adique}.}

\deleted{Soit $X$ projective lisse connexe sur $k$ alg clos, et soit $z_1\in Z_1^{(\mathrm{num})}(X)$, défini par $Z_1\in\mathrm{Cyc}_1(X)$.}

\uline{Prop.} 2.1. Soient $X,Y$ projectives lisses connexes sur $k$, de dim $n$ resp. $m$, et \deleted{a} $Z_1\in\mathrm{Cyc}_1(X\times Y)$. Alors il existe un homomorphisme
% slanted
\margin{\unsure{echanger|échanger} le rôle de $X$ et $Y$ ?}
unique
\diagram{\begin{tikzcd}[ampersand replacement=\&,column sep=small,row sep=tiny]
u: \& B_X\arrow[r] \& A_Y\\
\& \big\Vert \& \\
\& \underline{\mathrm{Pic}}^{00}_{X/k} \&
\end{tikzcd}}
satisfaisant la condition suivante : pour tout diviseur $D$ sur $X$, algébriquement équiv. à $0$, tel que $ {}^{t}Z(D)=\deleted{\unsure{\pi}}\mathrm{pr}_{2*}((D\times Y)\cdot Z)$ soit défini, on \unsure{a|ait} $u(\mathrm{cl}(D))=S_{\unsure{Y|X}}({}^{t}Z(D))$. Cet homomorphisme
% archive sheet number: 57
